\begin{lemma}
Let \(\mathcal A\) be any class of base vertices.  The number of red/blue
tagged projected pairs generated from witnesses in \(\mathcal A\) is at most
twice the number of final-vertex closed pairs generated by the same witnesses:
\[
\left|
\noderef{TwoBiteProjectedPairsClosedInClass}(\omega,I,\mathcal A)
\right|
\le
2\left|
\noderef{TwoBiteClosedPairsInClass}(\omega,I,\mathcal A)
\right|.
\]
\end{lemma}
\begin{proof}
Write \(Q_{\mathcal A}\) for
\(\noderef{TwoBiteProjectedPairsClosedInClass}(\omega,I,\mathcal A)\) and
write \(F_{\mathcal A}\) for
\(\noderef{TwoBiteClosedPairsInClass}(\omega,I,\mathcal A)\).  We prove a
two-to-one domination of \(Q_{\mathcal A}\) by \(F_{\mathcal A}\).  Recall
from \(\noderef{TwoBiteBasePairs}\) and \(\noderef{TwoBiteFinalPairs}\) that
pairs are represented by ordered endpoint pairs with increasing index, so each
unordered pair has at most one representative in the corresponding pair set.

Take a tagged projected pair \(e\in Q_{\mathcal A}\).  By the definition of
\(\noderef{TwoBiteProjectedPairsClosedInClass}\), \(e\) is an element of the
projected-pair set and it has at least one witnessing class vertex.  If \(e\)
is red-tagged, write \(e=\operatorname{red}(q)\).  Then there is some
\(x\in\mathcal A\) such that
\[
q\in
\noderef{TwoBiteBasePairs}
\bigl(\pi_R(X_x(I))\bigr).
\]
Unpacking \(\noderef{TwoBiteBasePairs}\), write \(q=(r_1,r_2)\) with
\(r_1,r_2\in\pi_R(X_x(I))\) and \(r_1<r_2\).  Choose
\(u,v\in X_x(I)\) with \(\pi_R(u)=r_1\) and \(\pi_R(v)=r_2\).  Since
\(r_1\ne r_2\), we have \(u\ne v\).  Let \(a\) be the increasing-index
representative of the final pair \(\{u,v\}\).  Then
\(a\in\noderef{TwoBiteFinalPairs}(X_x(I))\).  Because \(x\in\mathcal A\), the
definition of the bi-union
\(\noderef{TwoBiteClosedPairsInClass}(\omega,I,\mathcal A)\) gives
\(a\in F_{\mathcal A}\).  Moreover, the tagged projected pair \(e\) is exactly
the red tag of the increasing-index representative of
\(\{\pi_R(u),\pi_R(v)\}\).  Choose one such final pair \(a\) and set
\(\phi(e)=a\).

The blue-tagged case is identical.  If \(e=\operatorname{blue}(q)\), choose
\(x\in\mathcal A\) and vertices \(u,v\in X_x(I)\) witnessing
\(q\in\noderef{TwoBiteBasePairs}(\pi_B(X_x(I)))\); the distinct blue endpoints
force \(u\ne v\), and the increasing-index final pair \(a=\{u,v\}\) belongs to
\(F_{\mathcal A}\).  Set \(\phi(e)=a\).  Thus every element of
\(Q_{\mathcal A}\) has been assigned a chosen final pair in \(F_{\mathcal A}\).
The choice is arbitrary when several class vertices or several final pairs
witness the same \(e\); only the chosen final pair is used below.

Now fix a represented final pair \(a=(u,v)\in F_{\mathcal A}\).  Consider any
\(e\in Q_{\mathcal A}\) with \(\phi(e)=a\).  If \(e\) is red-tagged, then the
construction of \(\phi\) says that \(e\) is the red tag of the projected pair
coming from this very final pair \(a\).  If \(\pi_R(u)=\pi_R(v)\), no red
projected pair is produced.  If \(\pi_R(u)\ne\pi_R(v)\), then
\(\noderef{TwoBiteBasePairs}\) contains exactly one increasing-index
representative of the unordered pair
\(\{\pi_R(u),\pi_R(v)\}\).  Hence there is at most one red-tagged element in
the fiber \(\phi^{-1}(a)\).  The same argument with \(\pi_B\) gives at most
one blue-tagged element in \(\phi^{-1}(a)\).  The two tags are disjoint
summands of \(\noderef{TwoBiteProjectionPairSet}\), so
\[
|\phi^{-1}(a)|\le 1+1=2.
\]
This fiber bound is unaffected by unions or by multiple possible witnesses:
even if \(e\) is realized by many class vertices, once \(\phi(e)=a\) has been
chosen, the tag of \(e\) must be one of the two projected tags produced by the
fixed final pair \(a\).

The fibers \(\phi^{-1}(a)\), \(a\in F_{\mathcal A}\), partition
\(Q_{\mathcal A}\).  Summing the fiber bound over all final pairs gives
\[
|Q_{\mathcal A}|
=\sum_{f\in F_{\mathcal A}}|\phi^{-1}(f)|
\le \sum_{f\in F_{\mathcal A}}2
=2|F_{\mathcal A}|.
\]
This is the claimed inequality.
\end{proof}
